Considera la funció
\[
f(x)=x^3-3x+2 .
\]

\begin{apartats}

\apartat[1,25]{1}
Troba els punts de la gràfica de $f$ en què la recta tangent és paral·lela a la recta
$y=9x-4$.

\begin{solucio}
Rectes paral·leles tenen el mateix pendent: cal $f'(x)=9$.\\
$f'(x)=3x^2-3=9\iff x^2=4\iff x=\pm2$.\\
$f(2)=8-6+2=4$ i $f(-2)=-8+6+2=0$: els punts són $\boxed{(2,4)}$ i $\boxed{(-2,0)}$.
\end{solucio}

\begin{tria}{condicions-tasca}
\itemtria{original}{0,75}{1,25}
Escriu les equacions de les rectes tangents en aquests punts.

\begin{solucio}
En $(2,4)$: $y-4=9(x-2)$, és a dir $y=9x-14$.\\
En $(-2,0)$: $y=9(x+2)$, és a dir $y=9x+18$.
\end{solucio}

\itemtria{pendents-possibles}{0,75}{1,25}
Hi ha algun punt de la gràfica de $f$ en què la recta tangent tingui pendent $-4$? I pendent
$-3$? Justifica les respostes i, si n'hi ha, escriu l'equació de la tangent.

\begin{solucio}
El pendent de la tangent en el punt d'abscissa $x$ és $f'(x)=3x^2-3$. Com que $3x^2\ge0$,
sempre és $f'(x)\ge-3$.\\
Pendent $-4$: \textbf{cap punt}, perquè $f'(x)$ no pot ser més petit que $-3$.\\
Pendent $-3$: només si $3x^2=0$, és a dir en $x=0$. El punt és $(0,2)$ i la tangent és $y=-3x+2$.
\end{solucio}
\end{tria}

\begin{nomesllarg}
\apartat{0,75}
Troba els punts de la gràfica de $f$ en què la recta tangent és horitzontal.

\begin{solucio}
Tangent horitzontal vol dir pendent $0$: $f'(x)=3x^2-3=0\iff x=\pm1$.\\
$f(1)=0$ i $f(-1)=4$: els punts són $(1,0)$ i $(-1,4)$.
\end{solucio}
\end{nomesllarg}

\end{apartats}
